\section{ADR-013: macOS on Apple Silicon (experimental)}
\textbf{Status:} accepted; implemented as an experimental platform (Future
work). Windows stays the reference platform and behaves as before.

\textbf{Context.} Many musicians use a Mac. Every layer of Cover\_Maker was
written for Windows: the shell kills the engine's process tree with a Job
Object and hides console windows, the runtime bootstrap (ADR-006) bundles
\texttt{uv.exe} and pins Windows downloads, the engine keeps its data in
\texttt{\%LOCALAPPDATA\%} and finds the GPU with \texttt{nvidia-smi}, and the
lock file only resolves for \texttt{win32}. The project has no Mac: macOS can
only be built and tested in CI, on GitHub's \texttt{macos-latest} runners
(Apple Silicon virtual machines). Intel Macs are out of scope, because
PyTorch publishes no x86-64 macOS wheels after 2.2.

\textbf{Process lifecycle.} Unix has no Job Objects. The engine starts as
the leader of a new process group (\texttt{process\_group(0)}), which its
children join: the separation worker, ffmpeg and yt-dlp. Stopping the engine
sends \texttt{SIGKILL} to the whole group, and so does dropping the guard,
so the shell's exit path is the same as on Windows. A crash or
\texttt{kill -9} of the shell cannot clean up, so the existing parent-PID
watchdog (\texttt{COVER\_MAKER\_PARENT\_PID}, ADR-005) now kills the
engine's children before it exits, outside Windows. The shell's OS code lives
in \texttt{src/platform/} (\texttt{win.rs}, \texttt{unix.rs}), with the same
functions on each platform: the process tree guard, hidden windows, the data
folder, free space (\texttt{statvfs}), the executable bit, and the names
\texttt{Scripts}/\texttt{bin} and \texttt{.exe}.

\textbf{Separation on Apple Silicon.} audio-separator 0.48 already selects
PyTorch's Metal backend (\texttt{mps}) when it is available on an arm64
processor, and the ONNX Runtime CoreML provider when it is installed.
Native float16 is verified upstream only for BS- and Mel-RoFormer on
\texttt{mps}; for other model families it logs a warning and keeps float32.
So the GPU ladder (ADR-002) works unchanged: the \texttt{gpu} device means
CUDA on Windows and Metal on a Mac. Two differences are handled in the worker:
\texttt{PYTORCH\_ENABLE\_MPS\_FALLBACK=1} runs operations that Metal lacks
on the CPU instead of failing, and a CPU worker reports MPS as unavailable,
since Metal has no equivalent of \texttt{CUDA\_VISIBLE\_DEVICES=-1}. The
engine detects the Apple GPU without importing torch: on
\texttt{darwin}/\texttt{arm64} it is always there, named after the chip
(\texttt{sysctl machdep.cpu.brand\_string}). It shares the unified memory,
so its ``VRAM'' for the hardware tiers is two thirds of the RAM, which is
about Metal's recommended working set. An 8\,GB Mac is in the 4\,GB tier, a
16\,GB Mac in the 8\,GB tier, and 18\,GB or more in the 12\,GB tier. An
out-of-memory error on \texttt{mps} contains ``out of memory'', so the ladder
steps down as on CUDA, and \texttt{torch.mps.empty\_cache()} frees memory
between models.

\textbf{Runtime bootstrap.} The bootstrap does the same on both platforms.
The uv lock now resolves for
\texttt{sys\_platform == 'darwin' and platform\_machine == 'arm64'} too, with
a third PyTorch extra, \texttt{mps}, that takes PyTorch from PyPI (its macOS
wheel includes Metal; NVIDIA publishes no CUDA build for macOS, so the
\texttt{cu128} index applies only on Windows). The Windows profiles
(\texttt{core}, \texttt{cpu}, \texttt{cu128}) resolve to exactly the same
packages as before. The exported \texttt{pylock} files cover both platforms,
and the installer keeps only the wheels for the machine. The shell's only
torch build on macOS is \texttt{mps}, which also runs on the CPU. The
\texttt{mps} profile downloads about 290\,MB of wheels. uv installs CPython
3.12 from python-build-standalone, which has native arm64 builds.
\texttt{deps.json} now pins downloads per platform (\texttt{windows-x64},
\texttt{macos-arm64}), by URL and SHA-256. On macOS these are uv 0.13.6
(\texttt{.tar.gz}, staged at build time by \texttt{prepare-runtime.sh}) and
deno 2.9.8 (\texttt{.zip}). Extracted tools are marked executable. Files
that the app downloads itself carry no quarantine attribute, so Gatekeeper
does not check them.

Two gaps remain. \emph{ffmpeg:} the Windows runtime installs an LGPL shared
build from BtbN, which publishes no macOS builds, and the macOS builds found
(evermeet, osxexperts, Martin Riedl's build server) are GPL, with x264 and
x265. The project therefore builds its own (decided by the owner):
\texttt{desktop/scripts/build-ffmpeg-macos.sh} compiles FFmpeg n9.0.2 (pinned
by its release tag object) with LAME 3.100 for MP3 export (pinned by
SHA-256). Both are LGPL. It builds static programs with
\texttt{-{}-disable-autodetect}, linking only to macOS system libraries
(zlib, AudioToolbox), and the script refuses a build that is not LGPL or links
anything else. The zip has the layout of the Windows build (\texttt{bin/},
licence texts) plus \texttt{SOURCE.md} with the sources and configure line.
The \emph{ffmpeg for macOS} workflow builds it on \texttt{macos-latest},
checks that it encodes MP3, and keeps the zip with its SHA-256. The build is
not byte-for-byte reproducible, so one artifact is published once as a release
asset: the prerelease \texttt{ffmpeg-macos-n9.0.2} holds the zip of
workflow run 43429655833, and \texttt{deps.json} pins its SHA-256
(\texttt{230e18fc\ldots}). Unlike the Windows mirror, the URL does not follow
the app version, since the file never changes. A newer FFmpeg means a new
build, a new release tag and a new pin. The shell also appends Homebrew's
folders to the engine's \texttt{PATH}, since apps started from Finder get a
minimal one. \emph{diffq:} audio-separator
needs \texttt{diffq} outside Windows, and it ships only a source
distribution. The separation runtime therefore compiles it, which needs the
Xcode Command Line Tools. Homebrew installs them anyway. Mirroring a prebuilt
wheel would remove this step.

\textbf{Paths.} The engine (\texttt{paths.data\_root}), the shell
(\texttt{platform::local\_data\_dir}), the VST3 plug-in and the REAPER script
all use \texttt{\textasciitilde/Library/Application Support/Cover\_Maker}
(\texttt{Cover\_Maker.dev} for development builds) on macOS, with the same
layout as on Windows. The app bundle keeps the bootstrap files in
\texttt{Contents/Resources/bundle}.

\textbf{DAW integration.} The nih-plug VST3 plug-in builds on macOS as a
\texttt{.vst3} bundle, which CI checks. Its engine discovery and \emph{Show in
folder} (Finder) work there. Dragging stems out of the plug-in uses OLE and
stays Windows-only. Audio Unit support is not cheap: nih-plug has no AU
wrapper, so Logic Pro would need a separate wrapper (for example
clap-wrapper's AUv2 target), and it is left out. The REAPER script reads the
macOS discovery file and calls \texttt{/usr/bin/curl}; it has not been tested
on a Mac.

\textbf{Packaging.} Tauri builds the \texttt{.app} and a \texttt{.dmg}
(\texttt{--bundles app,dmg}). The minimum system is macOS 12.3, the first
with PyTorch's Metal backend. A manual run of the release workflow
(\texttt{workflow\_dispatch}) builds the disk image, the VST3 bundle and the
REAPER script, with their SHA-256 checksums, as a workflow artifact. It
attaches nothing to a release, and tag builds stay Windows-only. The app
updater has no macOS entry in \texttt{latest.json}, so \emph{Check for
updates} fails on macOS for now.

\textbf{Signing and notarization.} Signing for distribution needs an Apple
Developer Program membership (99\,USD a year, the owner's decision), a
Developer ID certificate, the hardened runtime and notarization with
\texttt{notarytool}. The bundled \texttt{uv} must be re-signed as nested
code. Without it the app is signed ad hoc (\texttt{signingIdentity "-"}),
which Apple Silicon requires to run code at all. Gatekeeper blocks a
downloaded copy (``cannot be opened because Apple cannot check it'').
The workaround, documented in the Installation Guide, is
\texttt{xattr -dr com.apple.quarantine /Applications/Cover\_Maker.app}, or
\emph{System Settings} $\rightarrow$ \emph{Privacy \& Security} $\rightarrow$
\emph{Open Anyway} on macOS 15. The same applies to the VST3 bundle.

\textbf{What CI verifies.} Every CI job now also runs on
\texttt{macos-latest}. The engine tests and the macOS separation runtime install
from the lock, which compiles diffq. PyTorch and audio-separator import there,
and the Apple GPU is detected. The UI checks and the build also run, as do the
shell's clippy and tests (process group, \texttt{statvfs}, paths), and the
VST3 plug-in's tests and bundle. The manual release run builds the
\texttt{.dmg}.

\textbf{What CI cannot verify.} The runners are virtual machines, which may
expose no Metal GPU, and no model is downloaded in CI, so no stems are
separated on \texttt{mps}. Its speed, memory use and fp16 quality are
unmeasured. Nobody has launched the app, run the first-run bootstrap
end to end or checked Gatekeeper's behaviour on a downloaded copy. The VST3
has not been loaded in a DAW, and the REAPER script has not been run. These
need a real Mac and are listed in the Installation Guide as experimental.

\textbf{Consequences.} One code base builds on both platforms, and the
Windows behaviour and lock are unchanged. Each OS-specific piece has one
place, so that piece can later move to a stable macOS build. Before macOS
leaves the experimental label, it needs
signing and notarization, a macOS entry for the updater and a test pass on
real hardware.
